\chapter{AdaptiveBridson}
\label{ch:adaptive-levelset}

\newthought{Trace the boundary, not the bulk.} \sampler{AdaptiveBridson} is the
feedback-driven contour tracer: it hunts the surface where a target function of your
observables crosses a chosen constant --- a $95\%$~CL contour, an exclusion edge, a
relic-density band --- and spends its budget densifying \emph{around that surface}
instead of the whole space. The current design (finalized 2026-07-22, validated on real
\textsc{idm}/micrOMEGAs relic-density scans) is an \emph{outer/core Bridson densifier}:
a coarse global pass finds the contour cheaply, then a two-band classification
(\emph{core} vs \emph{frontier}) decides where to pack more points, at a spatial scale
that shrinks generation by generation until the local band around the best point is
thin enough. Public tuning is deliberately small --- two knobs beyond the target itself
--- with everything else defaulted from real-world validation runs.

\begin{jnote}
This chapter documents the \emph{current} algorithm. An earlier design (crossing-edge
midpoint refinement, keyed on \yamlkey{contour\_precision}/\yamlkey{function\_tolerance})
is retired as the proposal engine; old checkpoints and cards that still set those keys
keep working (Section~\ref{sec:als-legacy}), but new cards should use the interface
below.
\end{jnote}

\section{Configuration}
\label{sec:als-config}

\begin{yamlwin}
Sampling:
  Method: AdaptiveBridson
  Variables:                        # REQUIRED: 2 to 5 entries
    - name: x
      distribution: { type: Flat, parameters: { min: 0.0, max: 1.0 } }
    - name: y
      distribution: { type: Flat, parameters: { min: 0.0, max: 1.0 } }

  Bounds:
    target_expression: "r2"         # REQUIRED: shared-language expression
    target_value: 0.04              # REQUIRED: the level T
    Seed: 7
    # Optional public tuning (defaults are fine for most cards):
    # outer_half_width: 0.02        # discovery band |f-T| <= w_outer
    # min_radius: 0.005             # u-space resolution floor
\end{yamlwin}

Every key in this chapter's tables goes directly in \yamlkey{Bounds} --- there is no
nested block underneath it, and \sampler{AdaptiveBridson} accepts the block strictly, so
a mistyped tuning key is reported before the scan starts rather than silently taking its
default (Section~\ref{sec:bounds}).

\begin{keylist}
  \item[target\_expression] \req\ (\code{JV2-MTH-041}). A shared-language expression
        (Chapter~\ref{ch:expressions}) evaluated in the control process over the
        observables your workflow \emph{returns} --- likelihood terms included. Its
        value is $f(u)$.
  \item[target\_value] \req\ (\code{JV2-MTH-042}). The level $T$ you are tracing:
        $f(u) = T$.
  \item[outer\_half\_width] \opt, default \code{0.02}. The \emph{discovery} band ---
        any evaluated point with $\lvert f - T\rvert \le w_{\mathrm{outer}}$ is worth
        keeping around as context; only its inner half (below) actually drives
        densification.
  \item[min\_radius] \opt, default $1/200$. The final u-space resolution floor: the
        per-generation spatial scale never shrinks past this, even if the local band is
        still not thin enough. Clamped down to \yamlkey{initial\_radius} if you set it
        larger.
\end{keylist}

Everything the tracer needs beyond the target and these two knobs is derived or safely
defaulted (Section~\ref{sec:als-keys}); prefer the two-knob card above and only add more
keys when you know specifically why.

\section{The algorithm}
\label{sec:als-algorithm}

Geometry --- Bridson spacing, windows, neighbor distances --- is always in
\textbf{u-space}; the physical map (Chapter~\ref{ch:variables}) is used only to build
the Sample that gets submitted and to read back $f$. The tracer works through a fixed
generation loop; within one generation it can take many \emph{fill passes} before the
spatial scale actually shrinks.

\begin{marginfigure}
\centering
\smallcaps{gen-0 Bridson}\\[2pt]
$\big\downarrow$\\[2pt]
\smallcaps{classify core / frontier}\\[2pt]
$\big\downarrow$\\[2pt]
\smallcaps{same-scale fill} $\circlearrowleft$\\
{\footnotesize root-corr, endpoints, bridge, densify}\\[2pt]
$\big\downarrow$ \emph{fill quiet}\\[2pt]
\smallcaps{band thin \& covered?}\\[2pt]
yes $\to$ \textbf{converged} \quad no $\to$ shrink $r_g$, repeat
\caption[The generation loop.]{One generation: classify, fill at fixed scale until
quiet, then either stop or shrink the scale and classify again.}
\end{marginfigure}

\subsection{Notation}

\begin{description}[style=nextline, leftmargin=1.2em]
  \item[$r_g$] the current spatial scale --- the Bridson/window radius for this
        generation.
  \item[$w_{\mathrm{outer}}, w_{\mathrm{core}}$] the discovery and densify-center
        half-widths in $f$; $w_{\mathrm{core}} = w_{\mathrm{outer}}/8$ by default.
  \item[core / frontier] a point is \emph{core} if $\lvert f-T\rvert \le w_{\mathrm{core}}$,
        \emph{frontier} if it is farther but still inside $w_{\mathrm{outer}}$; only
        cores act as densify centers and bridge endpoints.
  \item[best] $\arg\min_i \lvert f_i - T\rvert$ among points with a finite $f$.
  \item[generation / fill\_pass] a \emph{generation} is one completed shrink of $r_g$;
        a \emph{fill\_pass} is one densify / root-correction / endpoint / bridge round
        at the \textbf{same} $r_g$ --- fill passes do not advance the generation
        counter.
\end{description}

\begin{jwarn}
Two different masks look similar and are not: the \textbf{core/frontier} split
(above) decides who may open a densify window or act as a bridge center. The
\textbf{$t_{\min}, t_{\max}$} gate below is a completely separate measurement used only
to decide when to stop or shrink --- never confuse "is this point a core" with "is the
local band thin."
\end{jwarn}

\subsection{Generation 0}

A global Bridson fill at spacing \yamlkey{initial\_radius} covers the whole unit cube
($d \le 4$); at $d = 5$ a Sobol (quasi-random) fill is used instead, since Bridson's
rejection sampling degrades in five dimensions. Every gen-0 point is submitted and
evaluated; a point whose workflow fails never becomes best, core, or a densify center
--- non-finite $f$ is simply excluded, not guessed at.

\subsection{Classify, then gate}

Every evaluated point is classified by its \emph{absolute} distance from the target:
core if $\lvert f-T\rvert \le w_{\mathrm{core}}$, frontier if it is farther but still
within $w_{\mathrm{outer}}$. An empty core after gen-0 is not a failure --- edges whose
endpoints straddle $T$ seed the first cores via root-correction (below).

Separately, the tracer computes an exit/next-generation gate from \textbf{best}: collect
every finite $f$ within a ball of radius $2\,r_g$ around best, and take
$t_{\min}=\min f$, $t_{\max}=\max f$ over that set. Using $2\,r_g$ (not $r_g$) matters
--- under strict blue-noise spacing $r_g$, a ball of that radius around best is
essentially a single point, so its width would always read zero; $2\,r_g$ actually
measures local variation at the current scale. Fewer than two finite values in that ball
means "no band yet," which never counts as converged.

\subsection{Same-scale fill}
\label{sec:als-fill}

While cores exist and fill is still making structural progress, the tracer runs
repeated fill passes \emph{at the same $r_g$} --- generation does not advance:

\begin{enumerate}
  \item \textbf{Root-correction.} On the neighbor graph (Delaunay for $d\le3$, kNN for
        $d\ge4$) plus a nearest-opposite-side query, classify every edge:
        \begin{center}
        \begin{tabular}{@{}p{0.08\linewidth}p{0.40\linewidth}p{0.42\linewidth}@{}}
          \toprule
          \textbf{Class} & \textbf{Meaning} & \textbf{Action} \\
          \midrule
          A & neither endpoint is core & root-correction toward $T$ \\
          B & exactly one endpoint is core & expand from the core along the edge \\
          C & both cores, u-gap is large & spatial bridge between them \\
          D & both cores, already close & no fill needed on this edge \\
          \bottomrule
        \end{tabular}
        \end{center}
        A class-A or -B edge gets a new point placed by linear secant interpolation in
        $f$, clipped to stay well inside the edge:
        \[
        u = u_i + \alpha\,(u_j-u_i), \qquad
        \alpha = \mathrm{clip}_{[0.05,\,0.95]}\!\left(\frac{T-f_i}{f_j-f_i}\right)
        \]
        (falling back to the midpoint if blue-noise spacing rejects the secant point).
        A new point near $T$ becomes a core on the \emph{next} classify --- existing
        points are never moved.
  \item \textbf{Active endpoints.} Disconnected or open pieces of the contour need
        directional growth, not just densification around what already exists. The
        tracer keeps a persistent set of \emph{active endpoints} (contour tips) at the
        current $r_g$; each pass probes many directions on the hypersphere around every
        active endpoint (omni probes, default $\ge16$). An endpoint retires once a
        farther core supersedes it, it reaches the cube boundary, or it stalls for
        \yamlkey{endpoint\_stall\_passes} passes without progress. Endpoints are rebuilt
        from scratch whenever $r_g$ shrinks --- stale tips from a coarser scale are
        never carried forward.
  \item \textbf{Gap bridge.} When \yamlkey{bridge\_gaps} is on (default), nearby cores
        that are not yet well connected get a spatial bridge along a minimum-spanning
        tree, up to \yamlkey{bridge\_span\_factor}~$\times\,r_g$ apart.
  \item \textbf{Local densify.} Around every densify center, new candidates are drawn
        inside a ball of radius $r_g$ with Bridson-style trials
        (\yamlkey{k\_ref} per center), rejecting anything closer than $r_g$ to an
        already-accepted or in-flight point --- the same blue-noise separation as
        generation 0, just local. Over-dense centers are thinned first so proposal
        budget does not pile up on redundant cores.
\end{enumerate}

A fill pass's mid-pass accepts only constrain \emph{exclusion} for the rest of that pass
--- they never open new densify windows until the next classify. A point is never
resubmitted twice at the same $u$, even if it later falls out of the core band.

\textbf{Fill continues} while any of: new cores keep appearing, the best
$\lvert f-T\rvert$ keeps improving, contour extent or core-to-core spacing keeps
improving, or actionable A/B/C edges or active endpoints remain. After
\yamlkey{quiet\_fill\_passes} (default 3) with none of that, the scale is treated as
\emph{quiet} even if some straddling edges persist --- a lingering straddle alone is not
a reason to keep filling forever.

\subsection{Convergence, shrink, and hard stops}

Once a scale is quiet, the tracer checks \textbf{coverage}: enough cores
(\yamlkey{min\_cores\_for\_coverage}, default 4), forming one connected component, with
every core-to-core nearest-neighbor gap within
\yamlkey{core\_spacing\_factor}~$\times\,r_g$ (default $2.0$). Until coverage holds, the
scale never shrinks --- the tracer keeps appending cores instead.

With fill quiet and coverage OK, the tracer stops (\code{stop\_reason = converged}) once
\emph{all} of the following hold: fill is not needed, the best point is itself a core,
coverage is satisfied, and $t_{\max}-t_{\min} < \tau$ on the $2\,r_g$ ball about best
($\tau$ defaults to $w_{\mathrm{core}}$). If \yamlkey{final\_half\_width} is set, at
least one point must also land inside that tighter band.

Otherwise the tracer \emph{solidifies} the current cores into a permanent set, shrinks
$r_g \leftarrow \max(r_{\min},\, r_g \times \yamlkey{refinement\_factor})$ (default
$\times 0.5$), increments \yamlkey{generation}, resets \yamlkey{fill\_pass} to zero, and
rebuilds active endpoints from the new, finer geometry. If $r_g$ is already at
\yamlkey{min\_radius} and the band is still not thin, the tracer stops anyway --- it
never invents resolution finer than the floor you configured.

Convergence is not the only way a run ends. Table~\ref{tab:als-stops} lists the hard
stops, which are checked regardless of how the band is doing --- a run that ends on one
of these reports it in \code{stop\_reason}, so you can always tell a converged trace from
a truncated one.

\begin{table}[ht]
  \footnotesize
  \begingroup
  \setlength{\tabcolsep}{0pt}
  \begin{tabular}{@{}p{0.34\linewidth}p{0.66\linewidth}@{}}
    \toprule
    \textbf{Condition} & \textbf{\code{stop\_reason}} \\
    \midrule
    generation $\ge$ \yamlkey{max\_generations} & \code{max\_generations} \\
    total points $\ge$ \yamlkey{max\_points} & \code{max\_points} \\
    fill\_pass $\ge$ \yamlkey{max\_fill\_passes} & scale force-closed, then re-gated \\
    no finite $f$ at all & hard failure \\
    local windows saturated, no candidates left & \code{no\_refinement\_candidates} \\
    \bottomrule
  \end{tabular}
  \endgroup
  \caption{Hard safety stops, independent of convergence.}
  \label{tab:als-stops}
\end{table}

\begin{jtip}
When cores exist but every nearby bracket has cleared, \yamlkey{outer\_half\_width}
itself anneals toward \yamlkey{core\_half\_width} (by
\yamlkey{outer\_shrink\_factor}, default $0.7$, per pass) --- discovery support tightens
automatically as the contour becomes well sampled. This never substitutes for the
$t_{\min}/t_{\max}$ exit gate; it only trims how much frontier the tracer keeps around.
\end{jtip}

\section{How it behaves in the runtime}
\label{sec:als-runtime}

Every generation --- and every fill pass within it --- submits a batch, waits at the
usual feedback barrier for every proposal's observables, and only then classifies again;
Workers get \yamlkey{publish\_feedback} enabled automatically for this method. The
checkpoint carries every evaluated point and its $f$, the solidified core set, $r_g$,
\yamlkey{generation}, \yamlkey{fill\_pass}, the seed, and any pending
\textsc{uuid}s --- resume continues the exact same scale and fill state rather than
restarting the trace. \textsc{rng} streams are fixed per \code{(generation, fill\_pass)}
and core ordering is fixed, so a resumed run reproduces the same proposals a completed
one would have made.

At the default \code{INFO} level the log carries the start line, every $r_g$ transition,
and the final summary; fill-pass band diagnostics, bracket classification, endpoint
tracking, and densify timing are \code{DEBUG} (Chapter~\ref{ch:logging}).

\section{Outputs}
\label{sec:als-outputs}

Beyond the normal archive (\file{DATABASE/\allowbreak{}samples.\allowbreak{}hdf5},
plus \file{samples.\allowbreak{}csv} via \cli{Jarvis convert},
Section~\ref{sec:cli-convert}), the tracer writes
\file{<task\_result\_dir>/\allowbreak{}levelset.\allowbreak{}json} --- algorithm id
\code{outer\_core\_root\_correction\_bridson}, dimension, target and band values, the
final $r_g$ / $t_{\min}$ / $t_{\max}$, point counts, generation and fill-pass counts, the
\code{converged}/\code{stop\_reason} verdict, chained polylines for $d=2$, and a crossing
cloud with slice projections for $d \ge 4$. The stock \JPlot\ scenes overlay
\file{levelset.json} directly on scan scatter plots (Chapter~\ref{ch:jarvisplot}).
\yamlkey{save: true} calculator artifacts still land under \file{SAMPLE/} exactly as for
any other method (Chapter~\ref{ch:calculators}).

Both what the tracer emphasises in that file and how it builds its geometry depend on the
dimension; Table~\ref{tab:als-dimpolicy} collects that policy.

\begin{table}[ht]
  \footnotesize
  \begingroup
  \setlength{\tabcolsep}{0pt}
  \begin{tabular}{@{}p{0.104\linewidth}p{0.146\linewidth}p{0.208\linewidth}p{0.208\linewidth}p{0.334\linewidth}@{}}
    \toprule
    $d$ & Gen-0 & Neighbor graph & Densify engine & Output emphasis \\
    \midrule
    2 & Bridson & Delaunay & local Bridson & polylines + crossing cloud \\
    3 & Bridson & Delaunay & local Bridson & crossing cloud \\
    4 & Bridson & \textsc{k}nn & local Bridson & cloud + slice projections \\
    5 & Sobol & \textsc{k}nn & local Bridson & same as $d=4$ \\
    \bottomrule
  \end{tabular}
  \endgroup
  \caption{Dimension policy. Outside $2 \le d \le 5$, configuration raises
  \code{ValueError} before any point is drawn.}
  \label{tab:als-dimpolicy}
\end{table}

\section{Full key reference}
\label{sec:als-keys}

Table~\ref{tab:als-keys} is the complete \yamlkey{Sampling.Bounds} surface for this
method.

\begin{table}[ht]
  \footnotesize
  \begingroup
  \setlength{\tabcolsep}{0pt}
  \begin{tabular}{@{}p{0.219\linewidth}p{0.104\linewidth}p{0.125\linewidth}p{0.552\linewidth}@{}}
    \toprule
    \textbf{Key} & \textbf{Req.} & \textbf{Default} & \textbf{Notes} \\
    \midrule
    \multicolumn{4}{@{}l}{\emph{Public (prefer these on new cards)}} \\
    \yamlkey{target\_expression} & \req & --- & shared-language expression over
      returned observables \\
    \yamlkey{target\_value} & \req & --- & the level $T$ \\
    \yamlkey{outer\_half\_width} & & \code{0.02} & discovery band $\lvert f-T\rvert \le
      w_{\mathrm{outer}}$ \\
    \yamlkey{min\_radius} & & $1/200$ & u-space Euclidean floor for $r_g$; clamped to
      \yamlkey{initial\_radius} \\
    \multicolumn{4}{@{}l}{\emph{Common optional}} \\
    \yamlkey{initial\_radius} & & \code{0.10} & $r_0$: gen-0 fill and first $r_g$ \\
    \yamlkey{refinement\_factor} & & \code{0.5} & $r_g$ shrink multiplier per generation \\
    \yamlkey{bridge\_gaps} & & \code{true} & \textsc{mst} reconnection between cores \\
    \yamlkey{bridge\_span\_factor} & & \code{2.5} & max bridge length $= $ factor
      $\times\, r_g$ \\
    \yamlkey{max\_generations} & & \code{16} & max $r_g$ shrinks (scale index) \\
    \yamlkey{max\_points} & & \code{50000} & hard sample budget \\
    \yamlkey{max\_new\_per\_generation} & & \code{4000} & densify budget per fill pass \\
    \yamlkey{k\_ref} & & \code{30} & Bridson trials per densify center \\
    \yamlkey{neighbor\_graph} & & \code{auto} & \code{auto} $\to$ Delaunay ($d\le3$) /
      \textsc{k}nn ($d\ge4$); \code{delaunay}/\code{knn} force it \\
    \yamlkey{knn\_k} & & $4d$ & degree when the graph is \textsc{k}nn \\
    \yamlkey{Seed} & & \code{0} & alias \yamlkey{seed}; seeds every generation's draws \\
    \yamlkey{MaxWorker} & & workers & caps in-flight proposals \\
    \multicolumn{4}{@{}l}{\emph{Advanced / legacy (rarely needed on new cards)}} \\
    \yamlkey{core\_half\_width} & & $w_{\mathrm{outer}}/8$ & densify-center half-width;
      overrides the derived value \\
    \yamlkey{threshold} & & $= w_{\mathrm{core}}$ & the $t_{\max}-t_{\min}$ stop tolerance
      $\tau$; alias \yamlkey{function\_tolerance} \\
    \yamlkey{final\_half\_width} & & off & optional tighter export band $\le w_{\mathrm{core}}$ \\
    \yamlkey{radius\_shrink\_mode} & & \code{on\_coverage} & or legacy
      \code{every\_generation} \\
    \yamlkey{outer\_shrink\_factor} & & \code{0.7} & anneals $w_{\mathrm{outer}}$ once
      local brackets clear \\
    \yamlkey{core\_spacing\_factor} & & \code{2.0} & coverage: max core \textsc{nn} gap
      $\le$ factor $\times\, r_g$ \\
    \yamlkey{min\_cores\_for\_coverage} & & \code{4} & coverage floor \\
    \yamlkey{quiet\_fill\_passes} & & \code{3} & no-progress passes before a scale is quiet \\
    \yamlkey{max\_fill\_passes} & & $\max(64, 4g)$ & $g=$\,\yamlkey{max\_generations};
      safety cap on fill passes per scale \\
    \yamlkey{endpoint\_stall\_passes} & & \code{12} & passes before retiring an active
      endpoint \\
    \yamlkey{endpoint\_omni\_probes} & & $\max(16,\,4d)$ & directions probed per endpoint
      per pass \\
    \yamlkey{slice\_pairs} & & all pairs & $d\ge4$: which 2-D projections
      \file{levelset.json} exports \\
    \bottomrule
  \end{tabular}
  \endgroup
  \caption{\yamlkey{Sampling.Bounds} keys. \yamlkey{contour\_precision} (an
  older name for \yamlkey{threshold}) remains readable on legacy cards and checkpoints
  (Section~\ref{sec:als-legacy}).}
  \label{tab:als-keys}
\end{table}

\begin{jwarn}
\yamlkey{target\_expression} is compiled in the control process over the observables
your workflow \emph{returns}. Points whose workflow fails contribute nothing to any
graph or band --- a region where the calculator always crashes looks to the tracer like
missing territory, not evidence about $f$.
\end{jwarn}

\subsection{Reading an old card}
\label{sec:als-legacy}

Cards and checkpoints written against the retired edge-midpoint design still load:
\yamlkey{contour\_precision} is accepted as an alias of \yamlkey{threshold}, and
\yamlkey{function\_tolerance} likewise. Neither drives the densify engine anymore ---
the outer/core classification and the $t_{\min}/t_{\max}$ gate do that now --- but
setting them still moves $\tau$, so an old card keeps behaving sensibly rather than
silently ignoring its tuning. New cards should set \yamlkey{outer\_half\_width} (which
derives \yamlkey{threshold} automatically) instead of touching these directly.

\section{Tuning advice}

\begin{itemize}
  \item \textbf{Start from the two knobs, nothing else.} \yamlkey{outer\_half\_width}
        controls how much of $f$-space counts as "near the contour" for discovery;
        \yamlkey{min\_radius} controls the finest u-space resolution you are willing to
        pay for. Everything else (\yamlkey{core\_half\_width}, $\tau$, budgets) is
        already tuned from real relic-density scans --- override only when you know
        specifically why.
  \item \textbf{\yamlkey{min\_radius} is a floor, not a target.} The tracer stops
        refining once the band is thin \emph{or} $r_g$ hits this floor, whichever comes
        first --- setting it very small does not guarantee a tighter contour if
        \yamlkey{outer\_half\_width}/\yamlkey{max\_generations} exhaust first.
  \item \textbf{Barriers like fast workflows.} Every fill pass --- not just every
        generation --- waits for its slowest Sample. With hour-long calculators, a
        higher \yamlkey{max\_new\_per\_generation} means fewer, bigger barriers instead
        of many small ones.
  \item \textbf{Multiple contours mean multiple runs.} One \yamlkey{target\_value} per
        scan; the archived evaluations from one run are ordinary context you can replay
        into the next (\sampler{CSV}, Chapter~\ref{ch:csv-sampler}).
\end{itemize}

\section{Choosing \sampler{AdaptiveBridson} over the alternatives}

\begin{itemize}
  \item \textbf{vs \sampler{Bridson}/\sampler{Grid}/\sampler{Random}:} those spend their
        budget over the \emph{whole} declared space; \sampler{AdaptiveBridson} spends it
        on one target surface --- far cheaper for a contour, useless if you actually
        want the bulk.
  \item \textbf{vs the \textsc{mcmc} family / nested sampling
        (Chapter~\ref{ch:choosing-sampler}):} those characterize a \emph{posterior} (and,
        for nested sampling, its evidence); \sampler{AdaptiveBridson} traces a single
        level set of one expression --- not a posterior sampler, and it produces no
        evidence.
  \item \textbf{Dimension ceiling:} $2 \le d \le 5$ by construction --- past five
        dimensions a level set is no longer the right object to chase this way.
\end{itemize}

\section{Summary}
\label{sec:als-summary}

\sampler{AdaptiveBridson} is the only method here that reacts to results generation by
generation, spending its budget on a target surface instead of the whole space --- the
right choice for a contour or exclusion edge, never a general-purpose posterior sampler.
Table~\ref{tab:als-summary} is the short form.

\begin{table}[ht]
  \footnotesize
  \begingroup
  \setlength{\tabcolsep}{0pt}
  \begin{tabular}{@{}p{0.34\linewidth}p{0.66\linewidth}@{}}
    \toprule
    \textbf{Field} & \textbf{\sampler{AdaptiveBridson}} \\
    \midrule
    Kind & feedback \\
    Point count & until converged / generation \& point caps \\
    Dimensions & $2$--$5$ \\
    Required keys & \yamlkey{target\_expression}, \yamlkey{target\_value} \\
    Public tuning & \yamlkey{outer\_half\_width}, \yamlkey{min\_radius} \\
    u$\to$x mapping & yes \\
    \yamlkey{Seed} role & every generation's / fill pass's draws \\
    \yamlkey{selection} & filters proposals before submission \\
    Extra output & \file{levelset.json} \\
    \bottomrule
  \end{tabular}
  \endgroup
  \caption{\sampler{AdaptiveBridson} at a glance.}
  \label{tab:als-summary}
\end{table}
